\chapter*{Afterword}
\markboth{\textsc{AFTERWORD}}{}
\addcontentsline{toc}{chapter}{Afterword}

It is hard to believe, but we have finally reached the home stretch of
our \thepage-page long marathon and, preparing the final spurt, we
should first recall the main milestones that we have seen along this
way:
\begin{itemize}
\item As you might remember, we started off by summarizing the history
  of sentiment analysis, going back to its very origins in the ancient
  Greek philosophy and tracing its development to the present day;

\item Afterwards, to see what the current state of the art in opinion
  mining would yield on German Twitter, we created a corpus of
  $\approx8,000$ German tweets, collecting these messages for four
  different topics (federal elections, papal conclave, general
  political discussions, and casual everyday conversations).  To
  ensure a good recall of opinionated statements in the resulting
  dataset, we grouped all microblogs into three formal categories
  (tweets with a polar term from the SentiWS lexicon, messages
  containing a smiley, and all remaining microblogs) and sampled an
  equal number of tweets (666) for each of the four topics from each
  of these three categories.  After annotating the corpus in three
  steps (initial, adjudication, and final), we attained a reliable
  level of inter-annotator agreement for all elements (sentiments,
  sources, targets, polar terms, downtoners, negations, and
  intensifiers), finding that both selection criteria (topics and
  formal traits) significantly affected the distribution of sentiments
  and polar terms and the reliability of their annotation;

\item Then, at the first checkpoint, we compared existing German
  sentiment lexicons, which were translated from English resources and
  revised by human experts, with lexicons that were generated
  automatically from scratch with the help of state-of-the-art
  dictionary\mbox{-,} corpus\mbox{-,} and word-embedding--based
  methods.  An evaluation of these approaches on our corpus showed
  that semi-automatically translated polarity lists were generally
  better than the automatically induced ones, reaching 0.587
  macro-\F{} and attaining 0.955 micro-\F{}--score on the prediction
  of polar terms.  Furthermore, among fully automatic methods,
  dictionary-based systems showed stronger results than their corpus-
  and word-embedding--based competitors, yielding 0.479 macro-\F{} and
  0.962 micro-\F{}.  We could, however, improve on the latter metric
  (pushing it to 0.963) with our proposed linear projection solution,
  in which we first found a line that maximized the mutual distance
  between the projections of seed vectors with opposite semantic
  orientations and then projected the embeddings of all remaining
  words on that line, considering the distance of these projections to
  the median as polarity scores of the respective terms;

\item In Chapter \ref{chap:fgsa}, we turned our attention to the
  aspect-based sentiment analysis, in which we tried to predict the
  spans of sentiments, targets, and holders of opinions using two most
  popular approaches to this task---conditional random fields and
  recurrent neural networks.  We obtained our best results (0.287
  macro-\F{}) with the first-order linear-chain CRFs.  We could,
  however, increase these scores by using alternative topologies of
  CRFs (second-order linear-chain and semi-Markov CRFs) and also boost
  the macro-averaged \F{} to 0.38 by taking a narrower interpretation
  of sentiment spans (in which we only assigned the \textsc{Sentiment}
  tag to polar terms).  Further evaluation of these methods proved the
  utility of the text normalization step (which raised the macro-\F{}
  of the CRF-method by almost 3\%) and task-specific word embeddings
  with the least-squares fallback (which improved the
  macro-\F{}--score of the GRU system by 1.4\%);

\item Afterwards, in Chapter~\ref{chap:cgsa}, we addressed one of the
  most popular objective in contemporary sentiment
  analysis---message-level polarity prediction (which we called
  coarse-grained sentiment analysis [CGSA]).  To get a better overview
  of the numerous existing systems, we compared three larger families
  of CGSA methods---dictionary-, machine-learning--, and
  deep-learning--based ones, finding that the last two groups
  performed significantly better than the lexicon-based approaches
  (the best macro-\F{}--scores of machine- and deep-learning methods
  run up to 0.677 and 0.69 respectively, whereas the best
  lexicon-based solution [\citeauthor{Hu:04}, \citeyear{Hu:04}] only
  reached 0.641 macro-\F{}).  Apart from this, we improved the results
  of many reimplemented approaches by changing their default
  configuration (\eg{} abandoning polarity changing rules of
  lexicon-based systems, using alternative classifiers for ML-based
  classifiers, or taking the least-squares embeddings for DL-based
  methods).  In addition to the numerous reimplementations of popular
  existing algorithms, we also proposed our own
  solution---lexicon-based attention (LBA), in which we tried to unite
  the lexicon and deep-learning paradigms by taking a bidirectional
  LSTM network and explicitly pointing its attention to the polar
  terms that appeared in the analyzed messages.  With this solution,
  we not only outperformed all alternative DL systems, but also
  improved on the scores of ML-based classifiers, attaining 0.69
  macro-\F{} and 0.73 micro-\F{} on the PotTS corpus.  Similarly to
  our findings of the previous chapter, we observed a strong positive
  effect of text normalization and task-specific embeddings with the
  least-squares approximation;

\item Finally, in the last part, we tried to improve the results of
  the proposed LBA method by making it aware of the discourse
  structure.  For this purpose, we segmented all microblogs from the
  PotTS and SB10k corpora into elementary discourse units,
  individually analyzing each of these segments with our CGSA
  classifier, and then estimated the overall polarity of a tweet by
  joining the polarity scores of its EDUs over the RST tree.  We
  proposed three different ways of doing this joining (latent CRFs,
  latent-marginalized CRFs, and Recursive Dirichlet Process),
  obtaining better results than existing discourse-aware sentiment
  methods and also outperforming the original discourse-unaware
  baseline.  In the concluding experiments, we further improved these
  scores by using manually annotated RST trees and richer subsets of
  discourse relations.
\end{itemize}


\section*{Conclusions}
\addcontentsline{toc}{section}{Conclusions}

Now that we have gone past all these landmarks, it is time to unbag
the questions which we had asked ourselves at the beginning of this
endeavor, and try to answer them again, equipped with all knowledge
that we have acquired during our run.  Here we go:

\begin{itemize}
  \item\textbf{Can we apply opinion mining methods devised for
    standard English to German Twitter?}

    Yes, we can, but the success of these approaches might
    significantly vary depending on the task, the size and the
    reliability of the training data, as well as the evaluation metric
    that we use. For example, dictionary-based lexicon methods
    achieved fairly good results on their objective, but this success
    was mostly due to the high quality of the \textsc{GermaNet}
    annotation.  On the other hand, our manually labeled PotTS corpus
    was evidently too small for aspect-based sentiment systems, which
    failed to generalize to unseen tweets despite their very high
    scores on the training set.  Message-level sentiment approaches,
    vice versa, seemed to be quite happy with the size of the training
    dataset, attaining good results on both corpora (PotTS and SB10k).
    Nevertheless, we again experienced a lack of data while working on
    discourse-aware enhancements, many of which hit the same ceiling
    of the macro-averaged \F{}-scores.

    Apart from these difficulties arising from insufficient data, we
    also noticed a significant degradation of the scores for systems
    whose original tasks and evaluation metrics were different from
    ours.  For example, the lexicon generation method of
    \citet{Esuli:05} was originally designed to assign polarity scores
    to all \emph{synsets} found in the \textsc{WordNet} and not to
    produce a list of polar \emph{words}.  Similarly, the RNTN
    approach of \citet{Socher:13} was trained and evaluated on all
    syntactic subtrees of a document and not only at the top text
    level.  Likewise, the system of~\citet{Yessenalina:11} was devised
    for doing ordinal logistic regression and not polarity
    classification, as in our case.  As a result, all these approaches
    showed lower scores than their competitors in our evaluation, even
    though they are undoubtedly well suited for their original data
    and tasks.

    Due to the high diversity of methods, metrics, and tasks, it is
    difficult to provide a general recipe for transferring existing
    English sentiment systems to German Twitter, but we still would
    like to formulate at least a few rules of thumb, which came up
    during our experiments:
    \begin{itemize}
      \item\textbf{Prefer methods which are closest to your training
        objective} and which were trained under similar conditions
        w.r.t. the amount of data, their class distribution and
        domain;
      \item\textbf{Put every single setting of these methods into
        question}---bear in mind that things which work well in the
        original cases are not guaranteed to work in your
        situation.\footnote{In this respect, it is important to
          realize that every classification task is merely an attempt
          to solve a system of equations, so that methods which are
          good at solving one system might completely fail to solve
          other equations.}  The more options you try, the better will
        be your results;
      \item\textbf{Try using manually labeled resources for your
        target domain}, if they are available, but pay attention at
        the quality of their annotation---it often matters more than
        the corpus size;
      \item If there are manually annotated data, \textbf{prefer
        machine-learning methods to hard-coded rules}--- they will
        penalize their bad components automatically by themselves;
      \item\textbf{Do not use randomly initialized word embeddings for
        deep-learning systems}---initialize them with language-model
        vectors (which are cheap to obtain).  Otherwise, your model
        might get stuck in a very bad local optimum.
    \end{itemize}

  \item\textbf{Which groups of approaches are best suited for which
    sentiment tasks?}

    Based on our evaluation, we answer to this question as follows:
    \begin{itemize}
      \item\emph{Sentiment lexicon generation} is more amenable to
        dictionary-based solutions, provided that there exists a
        sufficiently big, reliably annotated lexical taxonomy for
        these systems.  If there is no such resource, one should
        better resort to word-embedding--based algorithms;

      \item With a limited amount of training data, \emph{aspect-based
        sentiment analysis} can be better addressed with probabilistic
        graphical models such as conditional random fields with
        hand-crafted features;

      \item On the other hand, plain \emph{message-level sentiment
        analysis} can be efficiently tackled with both machine- and
        deep-learning algorithms (such as SVM, logistic regression, or
        RNN);

      \item But probabilistic graphical models strike back at
        \emph{discourse-aware sentiment methods}, where they might
        even outperform pure neural-network solutions, although the
        margin of these improvements is not that large.
    \end{itemize}

    Thus, probabilistic model can still hold their ground when it
    comes to structured prediction, but the difference of these
    algorithms from and their improvements upon neural networks are
    gradually vanishing.

  \item\textbf{How much do word- and discourse-level analyses affect
    message-level sentiment classification?}

    Our evaluation in Section~\ref{cgsa:subsec:eval:lexicons} showed
    that the macro-averaged \F{}-scores of our proposed lexicon-based
    attention system varied by up to 14\% (from 0.64 to 0.69
    macro-\F{} on the PotTS corpus, and from 0.44 to 0.58 on SB10k)
    depending on the lexicon used by this approach.  At the same,
    discourse enhancements could only improve the results of LBA by at
    most 1.5\% percent (from 0.677 to 0.678 on PotTS, and from 0.557
    to 0.572 on SB10k).  Although it appears as if the lexicon
    component were more important to a sentiment system, we would like
    to preclude such incorrect conclusion, because
    \begin{inparaenum}[(a)]
      \item a full-fledged sentiment solution should take into account
        both linguistic levels (words and discourse) and
      \item these relative results might look different if we expand
        the analyzed domain to longer documents or apply
        discourse-aware methods to complete discussion threads.
    \end{inparaenum}

  \item\textbf{Does text normalization help analyze sentiments?}

    Yes, it definitely does.  As we could see in
    Chapters~\ref{chap:fgsa} and \ref{chap:cgsa}, normalization
    significantly improves the quality of aspect-based and
    message-level sentiment analyses, boosting the results on the
    former task by up to 4\% (see
    Table~\ref{snt-fgsa:tbl:normalization}) and improving the
    macro-averaged \F{}-measure of coarse-grained methods by up to
    25\% (see Table~\ref{snt-cgsa:tbl:res-no-normalization});

    The only question that remained unanswered in this context is
    which normalization steps exactly improve the scores of sentiment
    systems.  To make up for this omission, we separately deactivated
    each individual step of our text normalization pipeline
    (unification of Twitter phenomena, spelling correction, and
    normalization of slang terms) and rerun our message-level
    classification experiments using the lexicon-based attention
    system.  As we can see from the results in
    Table~\ref{afterword:tbl:lba-normalization-steps}, the
    micro-averaged \F{}-scores on both datasets benefit most from the
    unification of Twitter-specific phenomena, sinking by almost 19\%
    when this component is deactivated.  This step is also most useful
    for the macro-\F{} on the SB10k corpus, whereas the macro-average
    on PotTS mostly capitalizes on the normalization of slang terms.
    \begin{table}[htb!]
      \begin{center}
        \bgroup \setlength\tabcolsep{0.1\tabcolsep}\scriptsize
        \begin{tabular}{p{0.162\columnwidth} % first columm
            *{9}{>{\centering\arraybackslash}p{0.074\columnwidth}} % next nine columns
            *{2}{>{\centering\arraybackslash}p{0.068\columnwidth}}} % last two columns
          \toprule
          \multirow{2}*{\bfseries Method} & %
          \multicolumn{3}{c}{\bfseries Positive} & %
          \multicolumn{3}{c}{\bfseries Negative} & %
          \multicolumn{3}{c}{\bfseries Neutral} & %
          \multirow{2}{0.068\columnwidth}{\bfseries\centering Macro\newline \F{}$^{+/-}$} & %
          \multirow{2}{0.068\columnwidth}{\bfseries\centering Micro\newline \F{}}\\
          \cmidrule(lr){2-4}\cmidrule(lr){5-7}\cmidrule(lr){8-10}

          & Precision & Recall & \F{} & %
          Precision & Recall & \F{} & %
          Precision & Recall & \F{} & & \\\midrule

          \multicolumn{12}{c}{\cellcolor{cellcolor}PotTS}\\
          %% with normalization
          with normalization & 0.76 & 0.84 & 0.79 & %
          0.6 & 0.56 & 0.58 & %
          0.75 & 0.68 & 0.72 & %
          0.69 & 0.73\\
          %% without replacement of Twitter-specific phenomena
          %% General Statistics:
          %% precision    recall  f1-score   support
          %% positive       0.51      0.87      0.64       679
          %% negative       0.57      0.40      0.47       268
          %% neutral       0.68      0.22      0.34       593
          %% micro avg       0.54      0.54      0.54      1540
          %% macro avg       0.59      0.50      0.48      1540
          %% weighted avg       0.59      0.54      0.49      1540
          %% Macro-Averaged F1-Score (Positive and Negative Classes): 55.53%
          %% Micro-Averaged F1-Score (All Classes): 53.7662%
          w/o unification of Twitter phenomena & 0.51\negdelta{0.25} & 0.87\posdelta{0.03} & 0.64\negdelta{0.05} & %
          0.57\negdelta{0.03} & 0.4\negdelta{0.16} & 0.47\negdelta{0.11} & %
          0.68\negdelta{0.07} & 0.22\negdelta{0.46} & 0.34\negdelta{0.38} & %
          0.56\negdelta{0.13} & 0.54\negdelta{0.19}\\
          %% without spelling correction
          %% General Statistics:
          %% precision    recall  f1-score   support
          %% positive       0.67      0.84      0.74       648
          %% negative       0.61      0.34      0.44       268
          %% neutral       0.74      0.68      0.71       626
          %% micro avg       0.69      0.69      0.69      1542
          %% macro avg       0.67      0.62      0.63      1542
          %% weighted avg       0.69      0.69      0.68      1542
          %% Macro-Averaged F1-Score (Positive and Negative Classes): 59.01%
          %% Micro-Averaged F1-Score (All Classes): 68.9364%
          w/o spelling correction & 0.67\negdelta{0.09} & 0.84 & 0.74\negdelta{0.05} & %
          0.61\posdelta{0.01} & 0.34\negdelta{0.22} & 0.44\negdelta{0.14} & %
          0.74\negdelta{0.01} & 0.68 & 0.71\negdelta{0.01} & %
          0.59\negdelta{0.1} & 0.69\negdelta{0.04}\\
          %% without slang normalization
          %% General Statistics:
          %% precision    recall  f1-score   support
          %% positive       0.59      0.87      0.70       650
          %% negative       0.60      0.17      0.26       268
          %% neutral       0.72      0.60      0.65       624
          %% micro avg       0.64      0.64      0.64      1542
          %% macro avg       0.64      0.54      0.54      1542
          %% weighted avg       0.65      0.64      0.61      1542
          %% Macro-Averaged F1-Score (Positive and Negative Classes): 48.26%
          %% Micro-Averaged F1-Score (All Classes): 63.6187%
          w/o slang normalization & 0.59\negdelta{0.17} & 0.87\posdelta{0.03} & 0.7\negdelta{0.09} & %
          0.6 & 0.17\negdelta{0.39} & 0.26\negdelta{0.32} & %
          0.72\negdelta{0.03} & 0.6\negdelta{0.08} & 0.65\negdelta{0.07} & %
          0.48\negdelta{0.21} & 0.64\negdelta{0.09}\\

          \multicolumn{12}{c}{\cellcolor{cellcolor}SB10k}\\
          %% with normalization
          with normalization & 0.6 & 0.72 & 0.66 & %
          0.47 & 0.42 & 0.44 & %
          0.84 & 0.8 & 0.82 & %
          0.55 & 0.73\\
          %%%%%%%%%%%%%%%%%%%%%%%%%%%%%%%%%%%%%%%%%%%%%%%%%%%%%%%%%%%%%%%%%%
          %% without replacement of Twitter-specific phenomena
          %% cgsa_sentiment train -t lba --lstsq -l cgsa/data/lexicons/linproj.word2vec.kim_hovy_seedset.auto.txt ~/data/CGSA/SB10k/preprocessed-no-noise-cleaner/train/train.tsv  ~/data/CGSA/SB10k/preprocessed-no-noise-cleaner/dev/dev.tsv
          %% cgsa_sentiment test ~/data/CGSA/SB10k/preprocessed-no-noise-cleaner/test/test.tsv > ~/data/CGSA/SB10k/preprocessed-no-noise-cleaner/predicted/lba.test
          %% General Statistics:
          %% precision    recall  f1-score   support
          %% positive       0.36      0.85      0.50       354
          %% negative       0.60      0.25      0.35       212
          %% neutral       0.84      0.51      0.63       930
          %% micro avg       0.55      0.55      0.55      1496
          %% macro avg       0.60      0.53      0.49      1496
          %% weighted avg       0.69      0.55      0.56      1496
          %% Macro-Averaged F1-Score (Positive and Negative Classes): 42.52%
          %% Micro-Averaged F1-Score (All Classes): 55.1471%
          w/o unification of Twitter phenomena & 0.36\negdelta{0.24} & 0.85\posdelta{0.13} & 0.5\negdelta{0.16} &%
          0.6\posdelta{0.13} & 0.25\negdelta{0.17} & 0.35\negdelta{0.09} & %
          0.84 & 0.51\negdelta{0.29} & 0.63\negdelta{0.19} & %
          0.43\negdelta{0.12} & 0.55\negdelta{0.18}\\

          %%%%%%%%%%%%%%%%%%%%%%%%%%%%%%%%%%%%%%%%%%%%%%%%%%%%%%%%%%%%%%%%%%
          %% without spelling correction
          %% Commands:
          %% cgsa_sentiment train -t lba --lstsq -l cgsa/data/lexicons/linproj.word2vec.kim_hovy_seedset.auto.txt ~/data/CGSA/SB10k/preprocessed-no-misspelling-restorer/train/train.tsv  ~/data/CGSA/SB10k/preprocessed-no-misspelling-restorer/dev/dev.tsv
          %% cgsa_sentiment test ~/data/CGSA/SB10k/preprocessed-no-misspelling-restorer/test/test.tsv > ~/data/CGSA/SB10k/preprocessed-no-misspelling-restorer/predicted/lba.test
          %% General Statistics:
          %% precision    recall  f1-score   support
          %% positive       0.54      0.71      0.61       354
          %% negative       0.54      0.26      0.35       212
          %% neutral       0.79      0.79      0.79       930
          %% micro avg       0.70      0.70      0.70      1496
          %% macro avg       0.62      0.59      0.58      1496
          %% weighted avg       0.70      0.70      0.69      1496
          %% Macro-Averaged F1-Score (Positive and Negative Classes): 48.12%
          %% Micro-Averaged F1-Score (All Classes): 69.5856%
          w/o spelling correction & 0.54\negdelta{0.06} & 0.71\negdelta{0.01} & 0.61\negdelta{0.05} & %
          0.54\posdelta{0.07} & 0.26\negdelta{0.16} & 0.35\negdelta{0.09} & %
          0.79\negdelta{0.05} & 0.79\negdelta{0.01} & 0.79\negdelta{0.03} & %
          0.48\negdelta{0.07} & 0.7\negdelta{0.03}\\

          %%%%%%%%%%%%%%%%%%%%%%%%%%%%%%%%%%%%%%%%%%%%%%%%%%%%%%%%%%%%%%%%%%
          %% without slang normalization
          %% Commands:
          %% cgsa_sentiment train -t lba --lstsq -l cgsa/data/lexicons/linproj.word2vec.kim_hovy_seedset.auto.txt ~/data/CGSA/SB10k/preprocessed-no-slang-normalization/train/train.tsv  ~/data/CGSA/SB10k/preprocessed-no-slang-normalization/dev/dev.tsv
          %% cgsa_sentiment test ~/data/CGSA/SB10k/preprocessed-no-slang-normalization/test/test.tsv > ~/data/CGSA/SB10k/preprocessed-no-slang-normalization/predicted/lba.test
          %% General Statistics:
          %% precision    recall  f1-score   support
          %% positive       0.55      0.71      0.62       354
          %% negative       0.64      0.20      0.30       212
          %% neutral       0.78      0.82      0.80       930
          %% micro avg       0.70      0.70      0.70      1496
          %% macro avg       0.66      0.57      0.57      1496
          %% weighted avg       0.71      0.70      0.69      1496
          %% Macro-Averaged F1-Score (Positive and Negative Classes): 45.97%
          %% Micro-Averaged F1-Score (All Classes): 70.3877%
          w/o slang normalization & 0.55\negdelta{0.05} & 0.71\negdelta{0.01} & 0.62\negdelta{0.04} & %
          0.64\posdelta{0.17} & 0.2\negdelta{0.22} & 0.3\negdelta{0.14} & %
          0.78\negdelta{0.06} & 0.82\posdelta{0.02} & 0.8\negdelta{0.02} & %
          0.46\negdelta{0.09} & 0.7\negdelta{0.03}\\\bottomrule
        \end{tabular}
        \egroup
        \caption{LBA$^{(1)}$ results without single text normalization
          steps}
        \label{afterword:tbl:lba-normalization-steps}
      \end{center}
    \end{table}

  \item\textbf{Can we do better than existing approaches?}

    Yes, we can:
    \begin{itemize}
    \item we improved the macro-averaged results of exisitng
      lexicon-generation methods with our proposed linear projection
      algorithms;
    \item we increased the scores of aspect-based analysis by
      redefining the topologies of CRFs;
    \item our lexicon-based attention network outperformed many of
      its competitors on message-level classification;
    \item and, finally,we surpassed the discourse-unware baseline and
      other existing discourse-aware sentiment solutions with the
      proposed latent-marginalized CRFs and Recursive Dirichlet
      Process.
    \end{itemize}
\end{itemize}


\section*{Contributions}
\addcontentsline{toc}{section}{Contributions}

Apart from answering the above questions and pushing the state of the
art for several major sentiment tasks on the PotTS and SB10k corpora,
we have also paved the way for other researchers who want to work on
the same topics by releasing the data and the code that we used in our
experiments:
\begin{itemize}
\item the Potsdam Twitter Sentiment (PotTS) corpus is available at:\\
  \url{https://github.com/WladimirSidorenko/PotTS};
\item scripts and executables used in our lexicon generation chapter
  can be downloaded
  from:\\ \url{https://github.com/WladimirSidorenko/SentiLex};
\item for our text normalization pipeline and aspect-based sentiment
  methods, please refer to:\\
  \url{https://github.com/WladimirSidorenko/TextNormalization};
\item furthermore, you can find our coarse-grained sentiment
  approaches at:\\ \url{https://github.com/WladimirSidorenko/CGSA};
\item and get all discourse-aware solutions
  from:\\ \url{https://github.com/WladimirSidorenko/DASA};
\item last but not least, we have released the discourse segmenter,
  the adapted DPLP parser, and our modified version of
  \textsc{RST-Tool}, which was adjusted to the annotation of
  multilogues, at:
  \begin{itemize}
  \item\url{https://github.com/WladimirSidorenko/DiscourseSegmenter},
  \item\url{https://github.com/WladimirSidorenko/RSTParser}, and
  \item\url{https://github.com/WladimirSidorenko/RSTTool} respectively.
  \end{itemize}
\end{itemize}

In addition to open-sourcing all projects, we have also made a few
attempts to increase the visibility of our research with the following
publications:
\begin{itemize}
  \item the rule-based text normalization was described in
    \cite{Sidarenka:13};
  \item the PotTS corpus was presented in \cite{Sidarenka:16};
  \item in \cite{Sidarenka:16b}, we summarized the evaluation of
    existing sentiment lexicons (a separate paper on the linear
    projection algorithm was withdrawn by us due to a mistake in the
    initial implementation);
  \item in \cite{Sidarenka:16a} and \cite{Sidarenka:17}, we described
    our initial experiments on message-level classification;
  \item furthermore, we introduced the SVM-based discourse segmenter
    in \cite{Sidarenka:15};
  \item and sketched our pilot study of discourse annotation in
    Twitter in \cite{Sidarenka:15a}.
\end{itemize}
Unfortunately, due to the lack of experience at the initial stage of
working on this dissertation and because of the time shortage at the
concluding stage, I\footnote{Throughout this work we have been using
  the scientific ``we,'' considering the reader as a companion in our
  marathon.  But because at this point I start describing the
  limitations of this work, which I am the only person responsible
  for, I would like to exclude the reader from this criticism by
  switching to the solitary ``I''.} was not able to publish more or at
higher-level venues.  I apologize for that.


\section*{Limitations}
\addcontentsline{toc}{section}{Limitations}

Much to my regret, the initial lack of academic expertise has also
prevented me from running this scientific marathon faster, better,
and, most sadly, along more exciting places.  Alas, in the ``Sentiment
Analysis of German Twitter'', I have concentrated more on the
``sentiment'' part, much at the cost of ``German Twitter''.  I wish I
had tried out more sophisticated cross-lingual methods for adapting
English methods to German, elaborated more on linguistic traits of
microblogs, and addressed the social aspect of the Twitter network.
The main reason why I have not done all these things is that, six
years ago, whe I started working on this dissertation completely from
scratch, with neither data, nor code, nor any proper plan in my head,
I was so overwhelmed by the abundance of works on opinion mining and
text normalization that I decided to answer the questions whether
existing sentiment methods work, whether text normalization helps, and
whether I could improve on these methods, first; regrettably, these
questions have pretty much preoccupied me since then.  Another reason
why I refrained from addressing certain topics (why, for example, I
did not extend the Rhetorical Structure Theory to multilogues) is
because properly handling these problems would require another
dissertation and would certainly first need to be done for English in
order to get any international attention.


\section*{Final Remarks}
\addcontentsline{toc}{section}{Final Remarks}

Nevertheless, I believe that with this thesis I have basically built a
theme park on the moon.  Although this park is still missing a slot
machine and a few carousels, it does have a nice card table and many
other kinds of funny amusements.  Another good thing is that the new
carousels are now easier to build, the existing amusement rides work
for free, and the park itself might still entertain its occasional
guests.  You might have enjoyed this theme park as well, or you might
have not---I appreciate either of your choices and would like to thank
you in any case that you have visited it.
\begin{figure*}[htb]
  \centering \includegraphics[height=5em]{img/bender.png}
\end{figure*}
